下面是当前项目的操作手册 + 功能亮点 + 实现路径。按“用户怎么用、开发者怎么理解、后续怎么扩展”来梳理。

一、怎么打开

在项目目录：

powershell



cd F:\Project_Loverboy01
.\.venv312\Scripts\python.exe -m uvicorn app.api.server:app --host 127.0.0.1 --port 8000



然后浏览器打开：

text



http://127.0.0.1:8000/app/



也可以用启动脚本：

powershell



powershell -ExecutionPolicy Bypass -File .\scripts\start-agent-app.ps1



配置主要看：


.env.example (line 1)

.env.local-embeddings.example (line 1)


二、用户操作手册

页面分三大区：


Documents

用于上传和管理知识库文档。

支持：


直接粘贴文本

上传 TXT / Markdown / CSV / JSON / HTML / XML

上传 PDF

上传 DOCX


操作流程：

打开 Documents
点开 Add or Update Knowledge
填标题，或选择文件自动带出标题
可选填 metadata JSON
点击 Upload document
上传成功后会显示 chunk 数量和索引状态

上传后可以：


在 Stored Documents 里查看文档

在 Document Detail 里看全文

用 Search 做知识库检索

点 Reindex stored documents 重建索引




Agent Chat

用于和 Agent 对话。

常见问法：

text



请搜索知识库，总结这份文档的核心内容



text



请从我上传的 CCF 推荐目录里找出所有 C 类期刊，不要会议，输出完整列表



text



先查知识库里的部署步骤，再总结成 checklist



当前 Agent 支持：


会话记忆

多步工具调用

知识库搜索

通用全文抽取

计算器

高风险工具审批暂停




Paused Actions

用于处理高风险动作。

例如 send_email 工具是高风险工具。Agent 如果计划发邮件，会暂停，页面会出现审批卡片。

用户可以：


Approve

Reject

查看 run detail

查看技术事件流




三、功能亮点


本地优先的 Agent 应用

默认只允许本机访问 API。远程访问需要配置 token。

相关实现：


middleware.py (line 13)




文档上传能力比较完整

已支持文本、PDF、DOCX，并加了安全限制：


最大上传文件大小

最大抽取文本长度

PDF 最大页数


相关实现：


document_file_parser.py (line 55)

config.py (line 40)




RAG 检索不是纯向量，采用 hybrid

检索逻辑是：


lexical 关键词匹配

vector 相似度匹配

两者加权融合


相关实现：


retriever.py (line 9)




通用全文抽取能力

对“列出所有”“完整抽取”“排除某类”这类任务，不再只取少量 RAG chunk，而是直接扫全文。

例如：

text



找出所有 C 类期刊，不要会议



会走：

text



extract_document_items



相关实现：


document_extractor.py (line 65)




多步工具 Agent

Runtime 支持最多 MAX_TOOL_STEPS 轮：

text



plan -> tool -> plan -> tool -> answer



相关实现：


runtime.py (line 25)




审批幂等和并发保护

高风险工具审批时用原子 claim，重复点击不会重复执行。

相关实现：


approval_service.py (line 31)

run_repository.py (line 119)




工具系统可扩展

当前内置工具：


calculator

extract_document_items

extract_ccf_c_journals

search_knowledge_base

send_email


工具注册入口：


dependencies.py (line 104)




MCP HTTP 工具接入

当前支持：


mock MCP

HTTP / streamable_http MCP

tools/list discovery

tools/call


相关实现：


adapter.py (line 13)




四、整体实现路径

核心调用链是：

text



浏览器 /app
  -> FastAPI /api/chat/stream
  -> ChatService
  -> AgentRuntime
  -> LLMGateway plan
  -> ToolExecutor
  -> Built-in Tool / MCP Tool / RAG
  -> AgentRuntime answer
  -> SSE 返回前端
  -> Run + Message 持久化



对应代码：


前端页面：index.html (line 1)

前端逻辑：app.js (line 1)

API 入口：server.py (line 10)

聊天 API：chat.py (line 22)

聊天服务：chat_service.py (line 12)

Agent Runtime：runtime.py (line 14)

LLM Gateway：llm_gateway.py (line 39)

工具执行器：executor.py (line 8)


五、数据存储

当前主要是本地存储：


文档知识库

默认 SQLite：

text



data/knowledge_base.db



可切 Postgres。



Agent run / message

默认 SQLite：

text



data/multi_tool_agent.db





向量库

支持：


memory

qdrant_local

qdrant HTTP


当前推荐本地：

text



VECTOR_STORE_PROVIDER=qdrant_local





六、当前边界

已经适合本地或小团队试用，但还不是完整生产系统。

主要待补：


reindex job 状态仍是内存态，重启会丢

MCP stdio transport 还没实现

通用表格理解还偏规则化，复杂 PDF 表格需要增强

还没有用户体系/多租户权限

没有导出长抽取结果为 CSV/XLSX 的能力


七、开发者扩展方式

如果要加一个新工具：

在 app/tools/builtins/ 新建工具文件
定义 ToolDefinition
写 handler，返回 ToolExecutionResult
在 dependencies.py (line 104) 注册
如果要模型主动调用，更新 prompt 或 planner 规则
补测试